\chapter{National Accounts in Depth}\label{ch:b1:national-accounts-in-depth}

On 12 July 2016 Ireland's statistical office reported that the country's
real GDP had grown by $26.3\%$ in 2015; four months earlier it had put the
growth at $7.8\%$. % ledger: b1:ireland-2015
On 5 November 2010 Ghana's statistical service announced a GDP for 2010 of
44.8 billion cedi, against a previous estimate of 25.6 billion, and Ghana
moved from the low-income to the middle-income group overnight.
% ledger: b1:ghana-2010
Neither economy changed in a day. What changed was the estimate: in
Ireland, a few large firms moved assets that count in its accounts; in
Ghana, the accounts were rebuilt on a more recent base year. National
accounts are a measurement system, with rules about what counts and
estimates where data run out. The macroeconomics of the next chapters
uses their numbers, and this chapter shows how they are made.

\begin{openingchart}{Two estimates of the same growth: average yearly growth
of real GDP per person, 1990--2019, from the World Development Indicators
and from the Penn World Table, 171 countries, with the 45-degree
line.}{World Bank, World Development Indicators (NY.GDP.PCAP.KD); Penn World
Table 11.0 (rgdpna, pop)}
\begin{tikzpicture}
\begin{axis}[omscatter, xmin=-4, xmax=10, ymin=-4, ymax=10, xlabel={WDI growth (log points a year)},
  ylabel={PWT growth (log points a year)}, width=0.56\linewidth, height=0.5\linewidth]
  \addplot[only marks, mark=*, mark size=1.3pt, omDef, fill opacity=0.6] table[x=wdi, y=pwt, col sep=comma] {figdata/bachelor-1/18-national-accounts-in-depth/out/growth_gap.csv};
  \draw[gray, dashed] (axis cs:-4,-4) -- (axis cs:10,10);
\end{axis}
\end{tikzpicture}
\end{openingchart}

Two respected sources, measuring the same thing over the same 29
years (growth in log points a year, 100 times the change in the logarithm
divided by 29), mostly agree: the median gap is $0.13$ points a year and the
correlation $0.909$. But $16.4\%$ of the 171 countries differ by more than half
a point a year, and $7.6\%$ by more than a whole point. % figdata: growth_gap_stats
For Mali one source gives $1.495\%$ a year and the other $6.17\%$.
% figdata: growth_gap_largest
Over 29 years such a gap is the difference between an economy that
grew by half and one that grew sixfold.

\begin{recall}
GDP, value added, the production, income and expenditure approaches and
the identity $Y = C + I + G + X - M$ were met at school, as was the
conversion of incomes at purchasing power. Fixed-basket and chained price
indices are compared in \cref{ch:b1:consumer-welfare}; the deflator was met at school.
\end{recall}

\section{Sectors and the sequence of accounts}\label{sec:b1:national-accounts-in-depth:sectors}

\begin{definition}[Institutional sector]\label{def:b1:national-accounts-in-depth:sector}
\index{institutional sector}
An \emph{institutional sector} is a group of units that take economic
decisions of the same kind: households, non-financial corporations,
financial corporations, general government; units outside the country
form the \emph{rest of the world}.
\end{definition}

Each sector has the same \emph{sequence of accounts}. The production account
records its value added. The distribution accounts record the incomes it
receives and pays: wages, profits, interest, rents, taxes, social
contributions and benefits. The use-of-income account records what it
consumes and, as the balance, what it saves. The capital account records
its investment and, as the balance, what it lends to or borrows from the
other sectors.

\begin{definition}[Gross national income and disposable income]\label{def:b1:national-accounts-in-depth:gni}
\index{gross national income}\index{GNI}\index{disposable income}
\emph{Gross national income} is GDP plus the primary incomes (wages,
profits, interest) that residents receive from abroad, minus those paid to
non-residents. A sector's \emph{disposable income} is its primary income
plus the current transfers it receives minus those it pays; for the nation,
gross national disposable income is GNI plus net current transfers from
abroad.
\end{definition}

GDP measures what is produced in the country; \omterm{def:b1:national-accounts-in-depth:gni}{GNI} what its residents earn.
In Ireland they differ by a fifth or more, because much of the profit of
firms producing there belongs to non-residents: in 2015 Irish GNP grew by
$18.7\%$, against GDP's $26.3\%$. % ledger: b1:ireland-2015

\begin{omfigure}
\begin{tikzpicture}[x=1cm, y=0.012cm, font=\footnotesize]
  \fill[omDef!35] (0,0) rectangle (1.3,393.3);
  \fill[omThm!35] (1.6,387.4) rectangle (2.9,393.3);
  \fill[omDef!35] (3.2,0) rectangle (4.5,387.4);
  \fill[omProp!35] (4.8,387.4) rectangle (6.1,388.6);
  \fill[omDef!35] (6.4,0) rectangle (7.7,388.6);
  \fill[omThm!35] (8.0,84.9) rectangle (9.3,388.6);
  \fill[omDef!35] (9.6,0) rectangle (10.9,84.9);
  \node[above] at (0.65,393.3) {GDP 393.3};
  \node[above] at (2.25,393.3) {$-5.9$};
  \node[above] at (3.85,393.3) {GNI 387.4};
  \node[above] at (5.45,393.3) {$+1.2$};
  \node[above] at (7.05,393.3) {GNDI 388.6};
  \node[above] at (8.65,393.3) {$C + G$};
  \node[above] at (10.25,84.9) {saving 84.9};
  \node[below, text width=1.8cm, align=center, font=\scriptsize] at (2.25,0) {primary income abroad};
  \node[below, text width=1.8cm, align=center, font=\scriptsize] at (5.45,0) {transfers from abroad};
  \node[below, text width=1.8cm, align=center, font=\scriptsize] at (8.65,0) {consumption};
  \draw (-0.1,0) -- (11,0);
\end{tikzpicture}

{\small The sequence of accounts as a waterfall, Brenland 2025, billions of
dalens (the weekend problem): from output to national income, to
\omterm{def:b1:national-accounts-in-depth:gni}{disposable income}, to saving.}
\end{omfigure}

\section{Saving, investment and net lending}\label{sec:b1:national-accounts-in-depth:net-lending}

\begin{definition}[Net lending]\label{def:b1:national-accounts-in-depth:net-lending}
\index{net lending}\index{net borrowing}
A sector's \emph{net lending} is its saving minus its investment (plus net
capital transfers received): the amount it makes available to other
sectors. When negative it is \emph{net borrowing}.
\end{definition}

\begin{proposition}[The sector balances sum to zero]\label{prop:b1:national-accounts-in-depth:balances}
Ignoring capital transfers, the \omterm{def:b1:national-accounts-in-depth:net-lending}{net lending} of the domestic sectors adds up
to the current-account balance, and so to minus the \omterm{def:b1:national-accounts-in-depth:net-lending}{net lending} of the rest
of the world.
\end{proposition}

\begin{proof}
For the nation, gross national \omterm{def:b1:national-accounts-in-depth:gni}{disposable income} is
$Y + \mathrm{NPI} + \mathrm{TR} = C + I + G + X - M + \mathrm{NPI} + \mathrm{TR}$, and saving is that income less
$C + G$. So $S - I = X - M + \mathrm{NPI} + \mathrm{TR}$, the current account. National saving and
investment are the sums of the sectors', so the sectors' \omterm{def:b1:national-accounts-in-depth:net-lending}{net lending} sums
to the current account. What residents lend abroad, non-residents
borrow: the rest of the world's \omterm{def:b1:national-accounts-in-depth:net-lending}{net lending} is minus the current account.
\end{proof}

The proposition is an identity: it holds whatever happens. It says, for
example, that a government deficit must be matched by a private surplus,
a current-account deficit, or both. It does not say which, nor which
balance causes which.

\begin{omchart}{Germany: net lending by sector, per cent of GDP, 1999--2023;
the private sectors computed as the current account minus the
government's balance.}{Eurostat, Government deficit/surplus
(gov\_10dd\_edpt1, B9, S13); World Bank, World Development Indicators
(BN.CAB.XOKA.GD.ZS)}
\begin{tikzpicture}
\begin{axis}[omtime, xlabel={Year}, ylabel={Per cent of GDP}, xmin=1999, xmax=2023, ymin=-12, ymax=12,
  width=0.72\linewidth, height=0.42\linewidth, legend style={font=\footnotesize, legend columns=3, at={(0.5,-0.24)}, anchor=north}]
  \draw[gray, dotted] (axis cs:1999,0) -- (axis cs:2023,0);
  \addplot+[thick, no marks] table[x=year, y=gov, col sep=comma] {figdata/bachelor-1/18-national-accounts-in-depth/out/balances_de.csv};
  \addplot+[thick, no marks] table[x=year, y=private, col sep=comma] {figdata/bachelor-1/18-national-accounts-in-depth/out/balances_de.csv};
  \addplot+[thick, no marks] table[x=year, y=row, col sep=comma] {figdata/bachelor-1/18-national-accounts-in-depth/out/balances_de.csv};
  \legend{Government, Private sectors, Rest of the world}
\end{axis}
\end{tikzpicture}
\end{omchart}

In 1999 Germany's government borrowed $1.9\%$ of GDP and the rest of the world
lent it $1.43\%$. By 2019 the government lent $1.3\%$ and the private sectors
$6.57\%$, so the rest of the world borrowed $7.87\%$ of Germany's GDP. In 2020 the
government's balance swung to $-4.3\%$ and the private surplus rose to $10.65\%$.
% figdata: balances_de
The three lines always sum to zero.

\section{Volumes and prices}\label{sec:b1:national-accounts-in-depth:volumes}

Real GDP measures quantities at constant prices. Which prices? At the
prices of a fixed base year, a good whose price falls fast --- computers ---
keeps the high weight of the base year, and its growth inflates the total
more every year.

\begin{definition}[Chain-linked volume]\label{def:b1:national-accounts-in-depth:chain-linked}
\index{chain-linked volume}
A \emph{chain-linked volume} measures growth from each year to the next at
the previous year's prices, and multiplies these yearly growth factors to
link the years into one series, expressed in the prices of a reference year.
\end{definition}

\begin{measure}[Chain-linked volumes]
Two goods: computers, 10 at 10 dalens in 2015, 40 at 4 in 2020, 160 at 1.6 in
2025; everything else, 900 at 1, 950 at 1.1 and 1000 at 1.2. At fixed 2015
prices the volume goes 1000, 1350, 2600: growth of $160\%$. Chained, the volume
grows by $1350/1000$ from 2015 to 2020 (at 2015 prices) and by
$(160 \times 4 + 1000 \times 1.1)/(40 \times 4 + 950 \times 1.1) = 1740/1205$ from 2020 to 2025 (at 2020 prices), so it
reaches $1000 \times 1.35 \times 1.444 = 1949.4$: growth of $95\%$. The cost of chaining is
that components no longer add up: chained computers (1,600 in 2015
dalens) plus chained other goods (1,000) make 2,600, not 1,949.4.
% computed: code/bachelor-1/18-national-accounts-in-depth (problem, Part III)
\end{measure}

\begin{proposition}[Fixed weights overstate growth]\label{prop:b1:national-accounts-in-depth:fixed}
If quantities grow fastest for the goods whose relative prices fall, a
volume at fixed base-year prices grows faster than one at current-year
prices.
\end{proposition}

\begin{remark}
A chained volume reweights every period. When relative prices move the
same way in every period, as computer prices did, its growth lies between
the two fixed-weight measures; when they reverse, it need not.
\end{remark}

The proof is \cref{exo:b1:national-accounts-in-depth:fixed-base}; it is the
substitution bias of a fixed-basket price index, seen from the volume side.
The deflator is nominal GDP divided by the volume, so the same choice also
changes measured inflation.

\section{Purchasing-power parities}\label{sec:b1:national-accounts-in-depth:ppp}

To compare GDP across countries it must be put in one currency. Market
exchange rates price traded goods; most of GDP is services that do not
cross borders, and their prices differ far more.

\begin{definition}[PPP conversion factor]\label{def:b1:national-accounts-in-depth:ppp}
\index{PPP conversion factor}\index{price level index}
A \emph{PPP conversion factor}, or purchasing-power parity (PPP), between two
currencies is the ratio of the costs, each in its own currency, of the
same basket of goods and services in the two countries. The \emph{price level
index} of a country is its PPP divided by the market exchange rate.
\end{definition}

\begin{measure}[Building a PPP]
The International Comparison Program prices a common basket of goods and
services in each participating economy and combines the prices with each
economy's expenditure in its national accounts; its 2021 round covered 176
economies, the tenth comparison since 1968. % ledger: b1:icp-2021
A PPP is both a conversion factor and a spatial price index: it says how
many units of each currency buy what one unit of the base currency buys
in the base economy. Matching items across very different economies, and
pricing services such as government or housing, are its hardest parts.
\end{measure}

\begin{omchart}{Price level of household consumption (United States = 100)
against GDP per person at PPP, 199 countries, 2021.}{World Bank, World
Development Indicators (PA.NUS.PRVT.PLI, NY.GDP.PCAP.PP.KD)}
\begin{tikzpicture}
\begin{axis}[omscatter, xmode=log, log ticks with fixed point, xtickten={3,4,5}, xmin=700, xmax=160000, ymin=0, ymax=160,
  xlabel={GDP per person, PPP (log scale)}, ylabel={Price level (US = 100)}, width=0.7\linewidth, height=0.42\linewidth]
  \addplot[only marks, mark=*, mark size=1.3pt, omDef, fill opacity=0.6] table[x=gdp, y=pli, col sep=comma] {figdata/bachelor-1/18-national-accounts-in-depth/out/pli_2021.csv};
  \addplot[omThm, thick, sharp plot, mark=none] table[x=gdp, y=fit, col sep=comma] {figdata/bachelor-1/18-national-accounts-in-depth/out/pli_fit.csv};
  \addplot[only marks, mark=none, nodes near coords, point meta=explicit symbolic, nodes near coords style={font=\scriptsize, anchor=west}]
    table[x=gdp, y=pli, meta=country, col sep=comma] {figdata/bachelor-1/18-national-accounts-in-depth/out/pli_points.csv};
\end{axis}
\end{tikzpicture}
\end{omchart}

Richer countries are dearer: the price level rises by about 33 points for
each tenfold rise in income, and the correlation with log income is $0.58$.
% figdata: pli_stats
India's price level was $26.131$ against Switzerland's $136.199$. % figdata: pli_points
At market exchange rates India's GDP per person looks about four times
smaller relative to the United States than at PPP; why poorer countries are
cheaper is a question for \cref{ch:b1:open-economy}.

\section{What the accounts estimate}\label{sec:b1:national-accounts-in-depth:revisions}

\begin{definition}[Statistical revision]\label{def:b1:national-accounts-in-depth:revision}
\index{statistical revision}\index{benchmark revision}
A \emph{statistical revision} is a change to a published estimate as more
complete data arrive. A \emph{benchmark revision} (or rebasing) changes the
whole series at once: new sources, a new base year, new definitions.
\end{definition}

Quarterly estimates are revised as surveys and tax records come in; a
rebasing brings in parts of the economy the old sources missed, often the
informal sector and new industries. Ghana's rebasing moved the base year
from 1993 to 2006 and raised GDP per person in current dollars from about
550 to about 1,100. % ledger: b1:ghana-2010
Multinational firms create a third problem: where a firm records its
profits and its intellectual property decides in which country's GDP they
count.

\begin{record}[Ireland 2015 and Ghana 2010]
Ireland's GDP grew by $26.3\%$ in 2015 in the estimate of July 2016, against
$7.8\%$ in the estimate of March 2016; its GNP grew by $18.7\%$.
% ledger: b1:ireland-2015
Ghana's GDP was revised up by more than $60\%$ in November 2010 (Jerven and
Duncan 2012). % ledger: b1:ghana-2010
In both cases measured output jumped with no change in what anyone
produced or consumed that year: the accounts estimate an economy through
conventions, and changes in the conventions or in what they capture show
up as growth.
\end{record}

\begin{record}[How large are revisions?]
For the United States over 1993--2015, the average revision without
regard to sign, from the first quarterly estimates of GDP growth (at an
annual rate) to the latest ones, was between $1.21$ and $1.31$ percentage
points; from the advance to the third estimate, a few months later, it
was $0.65$ points (Fixler, Kanal and Tien 2018). % ledger: b1:fixler-revisions
A first estimate of $2\%$ growth is compatible with a final $1\%$ or $3\%$.
\end{record}

\section*{Exercises}\label{sec:b1:national-accounts-in-depth:exercises}

\begin{exercise}[$\star$]\label{exo:b1:national-accounts-in-depth:gni-ex}
GDP is 500 and residents' net primary income from abroad is $-20$. Find \omterm{def:b1:national-accounts-in-depth:gni}{GNI}.
\end{exercise}

\begin{exercise}[$\star$]\label{exo:b1:national-accounts-in-depth:three}
Households lend 30, firms borrow 10, government borrows 35 (net). Find the
current account and the rest of the world's \omterm{def:b1:national-accounts-in-depth:net-lending}{net lending}.
\end{exercise}

\begin{exercise}[$\star$]\label{exo:b1:national-accounts-in-depth:read}
From the opening chart's table, the largest gaps between the two sources
are Mali ($1.495\%$ in WDI, $6.17\%$ in PWT) and Honduras ($5.779\%$ against $1.236\%$).
% figdata: growth_gap_largest
Compute what each pair implies for GDP per person in 2019 relative to
1990.
\end{exercise}

\begin{exercise}[$\star$]\label{exo:b1:national-accounts-in-depth:recession}
Predict the sign of the change in the government's \omterm{def:b1:national-accounts-in-depth:net-lending}{net lending} in a
recession, and say which private balance moves the other way.
\end{exercise}

\begin{exercise}[$\star\star$]\label{exo:b1:national-accounts-in-depth:chain-ex}
Good A: 100 at 1, 110 at 1.1, 120 at 1.2 in years 0, 1, 2. Good B: 10 at 10, 20 at 6,
40 at 4. Compute the volume at year-0 prices and the chained volume, and show
that the chained components do not add up.
\end{exercise}

\begin{exercise}[$\star\star$]\label{exo:b1:national-accounts-in-depth:ppp-ex}
A basket holds 100 loaves, 10 haircuts and 1 phone. In Brenland they cost
2, 20 and 300 dalens; in Kestria 5, 30 and 900 kests. Compute the PPP. Kestria's
GDP is 60,000 million kests and the market rate is 3 kests per dalen: compute
its GDP in dalens both ways, and its price level relative to Brenland.
\end{exercise}

\begin{exercise}[$\star\star$]\label{exo:b1:national-accounts-in-depth:brenland}
From Brenland's accounts, nominal GDP was 379,562 million dalens in 2024 and
393,264 million in 2025, real GDP 286,895 and 292,390 million (2015 dalens), the
deflator 132.3 and 134.5. Compute \omterm{prop:b1:models-and-data:real-growth}{real growth}, deflator inflation and nominal
growth, and check the identity.
\end{exercise}

\begin{exercise}[$\star\star$]\label{exo:b1:national-accounts-in-depth:ip}
A multinational moves the patents behind its products to a subsidiary in
Brenland, which now books their profits. Predict the effect on Brenland's
GDP, \omterm{def:b1:national-accounts-in-depth:gni}{GNI} and employment.
\end{exercise}

\begin{exercise}[$\star\star$]\label{exo:b1:national-accounts-in-depth:pli-ex}
A country's PPP is 20 units of its currency per dollar and its market rate
50. Compute its \omterm{def:b1:national-accounts-in-depth:ppp}{price level index} and say what it means.
\end{exercise}

\begin{exercise}[$\star\star\star$]\label{exo:b1:national-accounts-in-depth:sum}
Prove that the \omterm{def:b1:national-accounts-in-depth:net-lending}{net lending} of households, firms, financial corporations
and government sums to the current account, with capital transfers
included.
\end{exercise}

\begin{exercise}[$\star\star\star$]\label{exo:b1:national-accounts-in-depth:fixed-base}
Prove \cref{prop:b1:national-accounts-in-depth:fixed} for two goods and two
years: if the good whose relative price falls has the faster quantity
growth, growth at base-year prices exceeds growth at current-year prices.
\end{exercise}

\begin{exercise}[$\star\star\star$]\label{exo:b1:national-accounts-in-depth:refute}
A country's measured growth is slow. What observations would show that
its growth is mismeasured rather than slow?
\end{exercise}

\section{Problem: Brenland's Statistician Rebases the Accounts}\label{sec:b1:national-accounts-in-depth:problem}

\begin{problem}[{Weekend problem --- a new level, the same growth}]\label{pb:b1:national-accounts-in-depth:rebase}
Brenland's accounts (invented) give, for 2025, in millions of dalens: GDP
393,264; $C = 225{,}014$, $I = 88{,}484$, $G = 78{,}653$, $X = 129{,}777$, $M = 128{,}664$; government
revenue 157,306 and spending 165,171; debt 236,818. Net primary income from
abroad is $-5{,}900$, consumption of fixed capital 59,000 and net current
transfers from abroad $+1{,}200$. In 2015 GDP was 248,438 (current and 2015
prices), and real GDP in 2025 was 292,390.

\medskip\noindent\textbf{Part I --- The sequence of accounts, 2025.}
\begin{enumerate}
  \item Check the expenditure identity.
  \item Compute \omterm{def:b1:national-accounts-in-depth:gni}{GNI}.
  \item Compute net national income.
  \item Compute gross national \omterm{def:b1:national-accounts-in-depth:gni}{disposable income}.
  \item Compute national saving.
  \item Compute the saving rate out of \omterm{def:b1:national-accounts-in-depth:gni}{disposable income}.
\end{enumerate}

\medskip\noindent\textbf{Part II --- Sector balances.}
\begin{enumerate}[resume]
  \item Compute the government's \omterm{def:b1:national-accounts-in-depth:net-lending}{net lending}, in millions and as a share of GDP.
  \item Compute the current account.
  \item Check that saving minus investment equals the current account.
  \item Compute the private sectors' \omterm{def:b1:national-accounts-in-depth:net-lending}{net lending}.
  \item Show that the private sectors, the government and the rest of the world sum to zero.
  \item Compute the debt ratio.
\end{enumerate}

\medskip\noindent\textbf{Part III --- Volumes.}
\begin{enumerate}[resume]
  \item By what factor did nominal GDP grow from 2015 to 2025?
  \item By what factor did real GDP grow, and at what average yearly rate?
  \item Check the 2025 deflator.
  \item Leave Brenland's GDP aside and take the toy two-good economy of the measure box on \omterm{def:b1:national-accounts-in-depth:chain-linked}{chain-linked volumes} (computers and other goods). Compute its volume at 2015 prices for 2015, 2020 and 2025.
  \item Compute the chained volume.
  \item Show that the chained components do not add up.
  \item Which measure should the statistician publish, and why?
\end{enumerate}

\medskip\noindent\textbf{Part IV --- Rebasing.} A new survey finds an informal
sector worth $12.8\%$ of the old real GDP in 2015 and $14\%$ in 2025.
\begin{enumerate}[resume]
  \item Compute the new real GDP for 2015 and 2025.
  \item By how much does the rebasing raise the 2025 level?
  \item Compute the new average yearly growth and compare.
  \item Compute the debt ratio after rebasing, assuming the same $14\%$ for nominal GDP.
  \item Why does the rebasing change the level so much more than the growth?
  \item State the result: the effect of the rebasing on the 2025 level and on growth from 2015 to 2025.
\end{enumerate}
\end{problem}
